\subsection{Reference Point Formation}
\label{subsec:reference-points}

%%%%%%%%%%%%%%%%%%%%%%%%%%%%%%%%%%%%%%%%%%%%%%%%%%%%%%%%%%%%%%%%%%%%%%%%%%%%%%%
% 10.5: How reference points are formed for each utility category
%%%%%%%%%%%%%%%%%%%%%%%%%%%%%%%%%%%%%%%%%%%%%%%%%%%%%%%%%%%%%%%%%%%%%%%%%%%%%%%

The uniform calculation logic (Section~\ref{subsec:calculation-logic}) established 
that all 144 components follow the same formula. The \emph{key difference} between 
INU, KNU, and IDN lies in how reference points are formed. This section develops 
the reference point formation process for each category.

\subsubsection{Why Reference Points Matter}

\paragraph{The Foundational Role.}
Reference points determine whether an outcome is experienced as Gain or Pain:

\begin{equation}
\text{Same outcome } x \quad \Rightarrow \quad 
\begin{cases}
\text{Gain if } x > r \\
\text{Pain if } x < r
\end{cases}
\end{equation}

A person earning \euro{70,000} experiences:
\begin{itemize}
\item \textbf{Gain} if their reference point is \euro{60,000}
\item \textbf{Pain} if their reference point is \euro{100,000}
\end{itemize}

\paragraph{Reference Points Are Category-Specific.}
The same person, evaluating the same outcome, may have different reference 
points for INU, KNU, and IDN---leading to simultaneous Gain in one category 
and Pain in another.

%%%%%%%%%%%%%%%%%%%%%%%%%%%%%%%%%%%%%%%%%%%%%%%%%%%%%%%%%%%%%%%%%%%%%%%%%%%%%%%
\subsubsection{INU Reference Points: Status Quo and Expectations}
%%%%%%%%%%%%%%%%%%%%%%%%%%%%%%%%%%%%%%%%%%%%%%%%%%%%%%%%%%%%%%%%%%%%%%%%%%%%%%%

\paragraph{The INU Reference Point Logic.}
Individual Utility evaluates outcomes against personal expectations:
``What did I have? What did I expect? What do my peers have?''

\begin{definition}[INU Reference Point]
\label{def:inu-reference}
\begin{equation}
r^{INU}_d = \alpha_1 \cdot SQ_d + \alpha_2 \cdot E_d + \alpha_3 \cdot SC_d
\end{equation}
where:
\begin{itemize}
\item $SQ_d$ = Status quo (current/recent state in dimension $d$)
\item $E_d$ = Expectations (anticipated future state)
\item $SC_d$ = Social comparison (peer group average)
\item $\alpha_1 + \alpha_2 + \alpha_3 = 1$ (weights)
\end{itemize}
\end{definition}

\paragraph{Component 1: Status Quo ($SQ_d$).}
The status quo is the default reference point---what one currently has or 
recently had:

\begin{itemize}
\item \textbf{Financial:} Current salary, wealth level
\item \textbf{Emotional:} Baseline mood, typical life satisfaction
\item \textbf{Physical:} Current health status
\item \textbf{Social:} Current relationship quality
\item \textbf{Digital:} Current connectivity level
\item \textbf{Ecological:} Current environmental quality
\end{itemize}

The status quo creates \emph{endowment effects}---losses from current state 
are felt more acutely than equivalent gains.

\paragraph{Component 2: Expectations ($E_d$).}
Expectations are forward-looking reference points---what one anticipated:

\begin{itemize}
\item \textbf{Formed by:} Past trends, promises, forecasts, plans
\item \textbf{Effect:} Outcomes below expectations feel like losses even if 
objectively positive
\item \textbf{Example:} Expected 5\% raise, got 3\% → Pain, even though income increased
\end{itemize}

\citet{koszegi2006model} formalized expectations-based reference points:
people feel loss when outcomes fall below what they expected.

\paragraph{Component 3: Social Comparison ($SC_d$).}
Social comparison evaluates outcomes relative to a peer group:

\begin{itemize}
\item \textbf{Formed by:} Observation of similar others, media exposure, social norms
\item \textbf{Effect:} Keeping up with the Joneses---relative position matters
\item \textbf{Example:} \euro{80k} salary feels like Pain if colleagues earn \euro{100k}
\end{itemize}

The empirical literature confirms that relative income predicts life satisfaction 
better than absolute income above subsistence \citep{clark2008relative}.

\paragraph{Typical INU Weights.}

\begin{table}[htbp]
\centering
\caption{INU Reference Point Weights by Context}
\label{tab:inu-weights}
\renewcommand{\arraystretch}{1.3}
\begin{tabular}{@{}llll@{}}
\toprule
\textbf{Context} & $\alpha_1$ (SQ) & $\alpha_2$ (Exp) & $\alpha_3$ (SC) \\
\midrule
Stable environment & 0.50 & 0.30 & 0.20 \\
High uncertainty & 0.60 & 0.20 & 0.20 \\
Competitive culture & 0.30 & 0.20 & 0.50 \\
Post-windfall & 0.40 & 0.40 & 0.20 \\
\bottomrule
\end{tabular}
\end{table}

%%%%%%%%%%%%%%%%%%%%%%%%%%%%%%%%%%%%%%%%%%%%%%%%%%%%%%%%%%%%%%%%%%%%%%%%%%%%%%%
\subsubsection{KNU Reference Points: Group Norms and Fairness}
%%%%%%%%%%%%%%%%%%%%%%%%%%%%%%%%%%%%%%%%%%%%%%%%%%%%%%%%%%%%%%%%%%%%%%%%%%%%%%%

\paragraph{The KNU Reference Point Logic.}
Collective Utility evaluates outcomes against group standards:
``What should we have? What is fair? How are similar groups doing?''

\begin{definition}[KNU Reference Point]
\label{def:knu-reference}
\begin{equation}
r^{KNU}_d = \beta_1 \cdot GN_d + \beta_2 \cdot FD_d + \beta_3 \cdot GC_d
\end{equation}
where:
\begin{itemize}
\item $GN_d$ = Group norm (what our group typically has/expects)
\item $FD_d$ = Fair distribution (equitable allocation standard)
\item $GC_d$ = Group comparison (what similar groups have)
\item $\beta_1 + \beta_2 + \beta_3 = 1$ (weights)
\end{itemize}
\end{definition}

\paragraph{Component 1: Group Norm ($GN_d$).}
The group norm is the collective expectation---what ``we'' should have:

\begin{itemize}
\item \textbf{Financial:} ``A family like ours needs \euro{X} to live decently''
\item \textbf{Emotional:} ``Our team should have good morale''
\item \textbf{Physical:} ``Our community should have access to healthcare''
\item \textbf{Social:} ``We should stick together''
\item \textbf{Digital:} ``Everyone in our village should have internet''
\item \textbf{Ecological:} ``Our region should have clean air''
\end{itemize}

\paragraph{Component 2: Fair Distribution ($FD_d$).}
Fairness is a KNU-specific reference point---not just what we have, but 
whether the distribution is just:

\begin{itemize}
\item \textbf{Equity:} Equal outcomes for equal inputs
\item \textbf{Equality:} Equal outcomes regardless of inputs
\item \textbf{Need:} Outcomes proportional to need
\end{itemize}

\citet{fehr1999theory} showed that people sacrifice personal payoffs to 
punish unfair distributions---even when they are not personally affected.

\paragraph{The Fairness Premium.}
KNU includes a fairness modifier:
\begin{equation}
r^{KNU}_d(\text{effective}) = r^{KNU}_d + \phi \cdot |\text{Actual distribution} - \text{Fair distribution}|
\end{equation}

Unfair distributions raise the effective reference point, making even 
objectively good outcomes feel like Pains.

\paragraph{Component 3: Group Comparison ($GC_d$).}
Groups compare themselves to other groups:

\begin{itemize}
\item \textbf{Between-group inequality:} ``Why does that region have better infrastructure?''
\item \textbf{Historical comparison:} ``We used to be the leading department''
\item \textbf{Aspirational comparison:} ``We should be like the top-performing teams''
\end{itemize}

\paragraph{Typical KNU Weights.}

\begin{table}[htbp]
\centering
\caption{KNU Reference Point Weights by Context}
\label{tab:knu-weights}
\renewcommand{\arraystretch}{1.3}
\begin{tabular}{@{}llll@{}}
\toprule
\textbf{Context} & $\beta_1$ (GN) & $\beta_2$ (FD) & $\beta_3$ (GC) \\
\midrule
Cohesive community & 0.50 & 0.30 & 0.20 \\
High inequality & 0.20 & 0.60 & 0.20 \\
Competitive groups & 0.30 & 0.20 & 0.50 \\
Traditional society & 0.60 & 0.25 & 0.15 \\
\bottomrule
\end{tabular}
\end{table}

%%%%%%%%%%%%%%%%%%%%%%%%%%%%%%%%%%%%%%%%%%%%%%%%%%%%%%%%%%%%%%%%%%%%%%%%%%%%%%%
\subsubsection{IDN Reference Points: Identity Standards}
%%%%%%%%%%%%%%%%%%%%%%%%%%%%%%%%%%%%%%%%%%%%%%%%%%%%%%%%%%%%%%%%%%%%%%%%%%%%%%%

\paragraph{The IDN Reference Point Logic.}
Identity Utility evaluates outcomes against self-concept:
``What should someone like me have/be? What does my identity prescribe?''

\begin{definition}[IDN Reference Point]
\label{def:idn-reference}
\begin{equation}
r^{IDN}_d = \gamma_1 \cdot IS_d + \gamma_2 \cdot RM_d + \gamma_3 \cdot NA_d
\end{equation}
where:
\begin{itemize}
\item $IS_d$ = Identity standard (what my identity prescribes)
\item $RM_d$ = Role model (what exemplars of my identity have)
\item $NA_d$ = Narrative anchor (what fits my life story)
\item $\gamma_1 + \gamma_2 + \gamma_3 = 1$ (weights)
\end{itemize}
\end{definition}

\paragraph{Component 1: Identity Standard ($IS_d$).}
The identity standard is the self-prescribed norm---what someone ``like me'' 
should have or be:

\begin{itemize}
\item \textbf{Financial:} ``A successful entrepreneur earns at least \euro{X}''
\item \textbf{Emotional:} ``A strong person doesn't show weakness''
\item \textbf{Physical:} ``An athlete maintains peak fitness''
\item \textbf{Social:} ``A good parent spends time with children''
\item \textbf{Digital:} ``A tech professional knows the latest tools''
\item \textbf{Ecological:} ``An environmentalist has a low carbon footprint''
\end{itemize}

\paragraph{Identity Standards Are Internalized Prescriptions.}
Following \citet{akerlofkranton}, identity standards come from:
\begin{itemize}
\item Social categories (gender, profession, nationality)
\item Chosen identities (vegan, runner, intellectual)
\item Ascribed identities (parent, child, citizen)
\end{itemize}

\paragraph{Component 2: Role Model ($RM_d$).}
Role models are concrete exemplars of the identity:

\begin{itemize}
\item \textbf{Function:} Provide specific, observable reference points
\item \textbf{Effect:} ``If X can do it, I should too''
\item \textbf{Danger:} Unrealistic role models create unattainable standards
\end{itemize}

\paragraph{Component 3: Narrative Anchor ($NA_d$).}
The narrative anchor is what fits one's life story:

\begin{itemize}
\item \textbf{Function:} Maintains coherence with past self and future aspirations
\item \textbf{Effect:} Changes that violate narrative feel identity-threatening
\item \textbf{Example:} A ``self-made'' person has high $r^{IDN}_F$ because their 
narrative requires continued success
\end{itemize}

\paragraph{Typical IDN Weights.}

\begin{table}[htbp]
\centering
\caption{IDN Reference Point Weights by Context}
\label{tab:idn-weights}
\renewcommand{\arraystretch}{1.3}
\begin{tabular}{@{}llll@{}}
\toprule
\textbf{Context} & $\gamma_1$ (IS) & $\gamma_2$ (RM) & $\gamma_3$ (NA) \\
\midrule
Stable identity & 0.50 & 0.25 & 0.25 \\
Identity transition & 0.30 & 0.40 & 0.30 \\
Strong role models & 0.30 & 0.50 & 0.20 \\
Narrative-focused & 0.25 & 0.25 & 0.50 \\
\bottomrule
\end{tabular}
\end{table}

%%%%%%%%%%%%%%%%%%%%%%%%%%%%%%%%%%%%%%%%%%%%%%%%%%%%%%%%%%%%%%%%%%%%%%%%%%%%%%%
\subsubsection{Reference Point Comparison: A Unified Example}
%%%%%%%%%%%%%%%%%%%%%%%%%%%%%%%%%%%%%%%%%%%%%%%%%%%%%%%%%%%%%%%%%%%%%%%%%%%%%%%

\paragraph{Scenario.}
Consider Maria, a 45-year-old marketing manager at a mid-sized company.
She earns \euro{75,000}. Let us trace her Financial reference points:

\begin{table}[htbp]
\centering
\caption{Maria's Financial Reference Points by Category}
\label{tab:maria-reference}
\renewcommand{\arraystretch}{1.4}
\begin{tabular}{@{}llll@{}}
\toprule
\textbf{Category} & \textbf{Components} & \textbf{Calculation} & $r_F$ \\
\midrule
\textbf{INU} & & & \\
\quad Status quo & Last year: \euro{72k} & $0.5 \times 72k$ & \\
\quad Expectations & Expected raise to \euro{78k} & $0.3 \times 78k$ & \\
\quad Social comparison & Peers earn \euro{80k} & $0.2 \times 80k$ & \\
\quad \textbf{Total} & & & \textbf{\euro{75.8k}} \\
\midrule
\textbf{KNU} & & & \\
\quad Group norm & Family needs \euro{65k} & $0.5 \times 65k$ & \\
\quad Fair distribution & Fair share is \euro{70k} & $0.3 \times 70k$ & \\
\quad Group comparison & Neighbor families: \euro{72k} & $0.2 \times 72k$ & \\
\quad \textbf{Total} & & & \textbf{\euro{67.9k}} \\
\midrule
\textbf{IDN} & & & \\
\quad Identity standard & ``Successful exec'': \euro{100k} & $0.5 \times 100k$ & \\
\quad Role model & Admired colleague: \euro{95k} & $0.3 \times 95k$ & \\
\quad Narrative & ``Rising career'': \euro{85k} & $0.2 \times 85k$ & \\
\quad \textbf{Total} & & & \textbf{\euro{95.5k}} \\
\bottomrule
\end{tabular}
\end{table}

\paragraph{Outcome Evaluation.}
Maria's actual salary: \euro{75,000}

\begin{table}[htbp]
\centering
\caption{Maria's Gains and Pains by Category}
\label{tab:maria-gains-pains}
\renewcommand{\arraystretch}{1.3}
\begin{tabular}{@{}llcccc@{}}
\toprule
\textbf{Category} & $r_F$ & $x_F$ & $G_F$ & $P_F$ & \textbf{Experience} \\
\midrule
INU & \euro{75.8k} & \euro{75k} & 0 & \euro{0.8k} & Slight Pain \\
KNU & \euro{67.9k} & \euro{75k} & \euro{7.1k} & 0 & Gain \\
IDN & \euro{95.5k} & \euro{75k} & 0 & \euro{20.5k} & \textbf{Strong Pain} \\
\bottomrule
\end{tabular}
\end{table}

\paragraph{Interpretation.}
The \emph{same} \euro{75,000} salary produces:
\begin{itemize}
\item Slight Pain for INU (below expectations and peers)
\item Gain for KNU (family is well-provided for)
\item \textbf{Strong Pain for IDN} (far below identity standard)
\end{itemize}

This explains why Maria might feel dissatisfied despite objectively good 
circumstances---her IDN reference point is driving negative utility.

%%%%%%%%%%%%%%%%%%%%%%%%%%%%%%%%%%%%%%%%%%%%%%%%%%%%%%%%%%%%%%%%%%%%%%%%%%%%%%%
\subsubsection{Reference Point Dynamics}
%%%%%%%%%%%%%%%%%%%%%%%%%%%%%%%%%%%%%%%%%%%%%%%%%%%%%%%%%%%%%%%%%%%%%%%%%%%%%%%

\paragraph{Reference Points Are Not Static.}
Reference points adapt over time:

\begin{equation}
r^i_{d,t+1} = r^i_{d,t} + \eta \cdot (x^i_{d,t} - r^i_{d,t})
\end{equation}

where $\eta \in (0, 1)$ is the adaptation rate.

\paragraph{Adaptation Rates Vary.}

\begin{table}[htbp]
\centering
\caption{Reference Point Adaptation Rates}
\label{tab:adaptation-rates}
\renewcommand{\arraystretch}{1.3}
\begin{tabular}{@{}lll@{}}
\toprule
\textbf{Category} & \textbf{Typical $\eta$} & \textbf{Interpretation} \\
\midrule
INU (Status quo) & 0.3--0.5 & Moderate adaptation to new normal \\
INU (Expectations) & 0.4--0.6 & Expectations update with experience \\
KNU (Group norm) & 0.2--0.4 & Group norms change slowly \\
IDN (Identity standard) & 0.1--0.3 & Identity standards are sticky \\
\bottomrule
\end{tabular}
\end{table}

\paragraph{Asymmetric Adaptation.}
Adaptation is often asymmetric:
\begin{itemize}
\item \textbf{Gains adapt faster:} The new car becomes ``normal'' quickly
\item \textbf{Losses adapt slower:} The pay cut still stings after years
\end{itemize}

This asymmetry explains the hedonic treadmill for gains but persistent 
unhappiness from losses.

\paragraph{Identity Reference Points Are Most Sticky.}
IDN reference points change most slowly because they are tied to self-concept:
\begin{itemize}
\item A ``successful person'' doesn't easily lower their standard
\item A ``coal miner'' doesn't easily adopt a ``coder'' identity standard
\item Narrative anchors require rewriting one's life story
\end{itemize}

This stickiness explains resistance to transitions that would improve 
INU but violate IDN.

%%%%%%%%%%%%%%%%%%%%%%%%%%%%%%%%%%%%%%%%%%%%%%%%%%%%%%%%%%%%%%%%%%%%%%%%%%%%%%%
\subsubsection{Context Effects on Reference Points}
%%%%%%%%%%%%%%%%%%%%%%%%%%%%%%%%%%%%%%%%%%%%%%%%%%%%%%%%%%%%%%%%%%%%%%%%%%%%%%%

\paragraph{Reference Points Are Context-Modulated.}
The $\Psi$-dimensions affect reference point formation:

\begin{equation}
r^i_d(\Psi) = \bar{r}^i_d + \sum_j \gamma^{r,i}_{dj} \cdot \Psi_j
\end{equation}

\paragraph{Key Context Effects.}

\begin{table}[htbp]
\centering
\caption{Context Effects on Reference Points}
\label{tab:context-reference}
\renewcommand{\arraystretch}{1.3}
\begin{tabular}{@{}lll@{}}
\toprule
\textbf{$\Psi$-Dimension} & \textbf{Effect on Reference Points} & \textbf{Mechanism} \\
\midrule
$\Psi_S$ (Social trust) & ↑ $r^{KNU}$ & Higher trust → higher collective expectations \\
$\Psi_I$ (Institutions) & ↓ $r^{INU}$ variability & Stable institutions → stable expectations \\
$\Psi_E$ (Economic) & ↑ $r^{INU}_F$, $r^{IDN}_F$ & Development raises financial aspirations \\
$\Psi_T$ (Temporal) & ↓ adaptation rate $\eta$ & Stability → slower reference adjustment \\
$\Psi_M$ (Market) & ↑ $r^{INU}_F$ via social comparison & Globalization expands comparison set \\
\bottomrule
\end{tabular}
\end{table}

\paragraph{Example: Globalization and Aspirations.}
Market integration ($\Psi_M$) raises INU financial reference points:
\begin{itemize}
\item Pre-globalization: Compare to local village → $r^{INU}_F$ moderate
\item Post-globalization: Compare to global middle class → $r^{INU}_F$ high
\item Result: Same objective income feels worse (larger Pain)
\end{itemize}

This helps explain rising discontent despite rising absolute incomes---the 
reference point rose faster than the outcome.

%%%%%%%%%%%%%%%%%%%%%%%%%%%%%%%%%%%%%%%%%%%%%%%%%%%%%%%%%%%%%%%%%%%%%%%%%%%%%%%
\subsubsection{Summary: Reference Point Formation}
%%%%%%%%%%%%%%%%%%%%%%%%%%%%%%%%%%%%%%%%%%%%%%%%%%%%%%%%%%%%%%%%%%%%%%%%%%%%%%%

\begin{tcolorbox}[colback=gray!5!white, colframe=gray!75!black, title=Section 10.5 Summary]
\textbf{Reference points determine whether outcomes are Gains or Pains.}

\textbf{Category-specific formation:}

\textbf{INU:} $r^{INU}_d = \alpha_1 \cdot SQ + \alpha_2 \cdot E + \alpha_3 \cdot SC$
\begin{itemize}
\item Status quo, Expectations, Social comparison
\item Question: ``What did I have/expect? What do peers have?''
\end{itemize}

\textbf{KNU:} $r^{KNU}_d = \beta_1 \cdot GN + \beta_2 \cdot FD + \beta_3 \cdot GC$
\begin{itemize}
\item Group norm, Fair distribution, Group comparison
\item Question: ``What should we have? Is it fair?''
\end{itemize}

\textbf{IDN:} $r^{IDN}_d = \gamma_1 \cdot IS + \gamma_2 \cdot RM + \gamma_3 \cdot NA$
\begin{itemize}
\item Identity standard, Role model, Narrative anchor
\item Question: ``What should someone like me have/be?''
\end{itemize}

\textbf{Key properties:}
\begin{itemize}
\item Same outcome → different reference points → different Gain/Pain
\item Reference points adapt over time (but IDN adapts slowest)
\item Context ($\Psi$) modulates reference point formation
\item IDN stickiness explains resistance to beneficial transitions
\end{itemize}
\end{tcolorbox}

\vspace{1em}
\hrule
\vspace{0.5em}
\noindent\textit{Next section: 10.6 Inter-Category Complementarities}

%%%%%%%%%%%%%%%%%%%%%%%%%%%%%%%%%%%%%%%%%%%%%%%%%%%%%%%%%%%%%%%%%%%%%%%%%%%%%%%
